\documentclass[12pt,a4paper]{article}
\usepackage{../latex/dsdlab-report}

\begin{document}
	\reporttitlepage{۲}{طراحی مدارهای ترتیبی با شماتیک}

	\section*{۲.۱ مقدمه و اهداف}
	هدف آزمایش، طراحی شماتیک یک مدار ترتیبی برای کنترل اتاق انتظار با ظرفیت ۱۵ نفر است. مدار باید عبور هم‌زمان از ورودی و خروجی، محدودیت ظرفیت، زمان مجاز ورود و کنترل درها را با شمارنده‌ی \lr{Up/Down} مدیریت کند.

	\section*{۲.۲ طراحی مدار و نمودار حالت}
	مدار در دو سطح طراحی شد. در سطح اول یک شمارنده‌ی چهار بیتی بالا/پایین با چهار فلیپ‌فلاپ \lr{JK} ساخته شد. ورودی \lr{U} جهت شمارش، \lr{Enable} اجازه‌ی شمارش، \lr{Clk} لبه‌ی فعال ساعت و \lr{Clr} پاک‌سازی شمارنده را تعیین می‌کنند. بیت کم‌ارزش در هر پالس مجاز تغییر می‌کند و شرط تغییر بیت‌های بالاتر از ترکیب بیت‌های کم‌ارزش‌تر و جهت شمارش به دست می‌آید. به این ترتیب، مدار در حالت \lr{U=1} صعودی و در حالت \lr{U=0} نزولی عمل می‌کند.

	\shot{screenshots/Counter_without_zoom.png}{شماتیک کامل شمارنده‌ی چهار بیتی بالا/پایین}
	\shot{screenshots/Counter_with_zoom_1.png}{جزئیات تولید ورودی فلیپ‌فلاپ‌های بیت صفر تا دو}
	\shot{screenshots/Counter_with_zoom_2.png}{جزئیات مرحله‌ی بیت پرارزش شمارنده}

	در سطح دوم شمارنده به‌صورت یک بلوک در مدار کنترل اتاق قرار گرفت. تنها عبور یک‌طرفه باعث تغییر تعداد افراد می‌شود؛ بنابراین برای ورود و خروج هم‌زمان، شمارنده بدون تغییر باقی می‌ماند. سیگنال \lr{NotFilled} با یک گیت \lr{NAND4} ساخته شده و تا پیش از رسیدن شمارنده به \(1111_2=15\) فعال است. همچنین \lr{Close} خروجی یک \lr{NOR4} است و دقیقاً در \(Q=0000\) فعال می‌شود تا در خروجی در حالت خالی بودن اتاق بسته بماند.

	درخواست باز شدن در ورودی از رابطه‌ی
	\[
		Request = Ent \land T \land NotFilled
	\]
	به دست می‌آید. خروجی \lr{Open} در یک فلیپ‌فلاپ \lr{D} نگه‌داری می‌شود تا درخواست لحظه‌ای دکمه به یک فرمان پایدار برای عبور فرد تبدیل شود. پس از مشاهده‌ی پالس حسگر \lr{IN}، مسیر بازخورد فرمان بازشدن را پاک می‌کند.

	\shot{screenshots/Citcut_without_zoom.png}{شماتیک کامل کنترل اتاق انتظار}
	\shot{screenshots/Circut_with_zoom_1.png}{بخش تولید اجازه‌ی شمارش و درخواست ورود}
	\shot{screenshots/Circut_with_zoom_2.png}{بلوک شمارنده، آشکارسازهای پر و خالی و حافظه‌ی فرمان بازشدن}

	نمای \lr{RTL} حاصل از سنتز نیز ساختار مورد انتظار را نشان می‌دهد: چهار عنصر حافظه و منطق ترکیبی در شمارنده، و بلوک شمارنده به‌همراه یک فلیپ‌فلاپ برای فرمان \lr{Open} در مدار اصلی.

	\shot{screenshots/RTL_view_Counter.png}{نمای \lr{RTL} شمارنده}
	\shot{screenshots/RTL_view_circut.png}{نمای \lr{RTL} مدار کامل}

	\section*{۲.۳ شبیه‌سازی و محاسبه‌ی فرکانس کاری}
	ابتدا شمارنده به‌صورت مستقل آزمایش شد. شکل موج نشان می‌دهد که با فعال بودن \lr{Enable} و \lr{U=1}، مقدار \(Q\) در لبه‌های ساعت به ترتیب افزایش می‌یابد؛ با تغییر \lr{U} مسیر شمارش معکوس می‌شود و با غیرفعال شدن \lr{Enable} مقدار قبلی حفظ می‌گردد. اعمال پاک‌سازی نیز خروجی را به صفر بازمی‌گرداند.

	\shot{screenshots/Counter_test.png}{شبیه‌سازی عملکرد شمارنده‌ی بالا/پایین}

	در شبیه‌سازی مدار کامل، ترکیب‌های مختلف \lr{IN}، \lr{OUT}، \lr{Ent} و \lr{T} اعمال شدند. خروجی \(Q\) در ورود خالص افزایش و در خروج خالص کاهش دارد. وقتی هر دو حسگر هم‌زمان فعال‌اند، تعداد افراد ثابت می‌ماند. فرمان \lr{Open} فقط با درخواست مجاز فعال می‌شود و \lr{Close} هنگام صفر شدن شمارنده یک می‌شود.

	\shot{screenshots/test_Circut.png}{شکل موج مدار کنترل اتاق برای سناریوهای ورود و خروج}

	برای تحلیل زمانی، در فایل \lr{SDC} یک ساعت با دوره‌ی \(10\,\mathrm{ns}\)، معادل \(100\,\mathrm{MHz}\)، تعریف و عدم‌قطعیت ساعت نیز استخراج شد.

	\shot{screenshots/file_for_clock.png}{قید ساعت در فایل \lr{Circut.sdc}}

	گزارش \lr{Fmax} مقدار محاسباتی \(446.83\,\mathrm{MHz}\) را نشان می‌دهد؛ بااین‌حال مقدار مجاز گزارش‌شده برای قطعه به‌علت محدودیت نرخ تغییر پایه‌های ورودی/خروجی \(250\,\mathrm{MHz}\) است. بنابراین فرکانس کاری قابل اتکا برای این پیاده‌سازی حداکثر \(250\,\mathrm{MHz}\) در نظر گرفته می‌شود و قید \(100\,\mathrm{MHz}\) انتخاب‌شده حاشیه‌ی کافی دارد.

	\shot{screenshots/فرکانس.png}{خلاصه‌ی بیشینه فرکانس مدار}

	\section*{۲.۴ سنتز و پیاده‌سازی نرم‌افزاری}
	هدف سنتز، قطعه‌ی \lr{EP4CGX15BF14C6} از خانواده‌ی \lr{Cyclone IV GX} بوده است. شمارنده با ۵ تابع ترکیبی و ۴ ثبات، و مدار کامل با ۱۱ تابع ترکیبی و ۵ ثبات پیاده‌سازی شده‌اند؛ مصرف منابع در هر دو حالت کمتر از یک درصد ظرفیت قطعه است. مدار کامل ۱۲ پایه و شمارنده‌ی مستقل ۸ پایه به کار می‌گیرند.

	\shot{screenshots/flow_summary_Counter.png}{خلاصه‌ی منابع مصرفی شمارنده}
	\shot{screenshots/flow_summary_Circut.png}{خلاصه‌ی منابع مصرفی مدار کامل}

	کامپایل هر دو طرح بدون خطا انجام شد. هشدارهای باقی‌مانده عمدتاً به کامل نبودن قیود \lr{setup/hold} مربوط‌اند و مانع تولید نت‌لیست نشده‌اند. با توجه به برگزاری مجازی آزمایشگاه، ارزیابی مدار مطابق روال درس بر پایه‌ی شبیه‌سازی، تحلیل زمانی و سنتز نرم‌افزاری انجام شد.

	\shot{screenshots/compile_success_counter.png}{کامپایل موفق شمارنده با صفر خطا}
	\shot{screenshots/compile_success_circut.png}{کامپایل موفق مدار کامل با صفر خطا}

	\section*{۲.۵ نتایج}
	\begin{itemize}
		\item شمارنده‌ی چهار بیتی در هر دو جهت و در حالت‌های پاک‌سازی و توقف مطابق انتظار عمل کرد.
		\item منطق ورود و خروج هم‌زمان از تغییر اشتباه تعداد افراد جلوگیری کرد.
		\item آشکارسازی ظرفیت ۱۵ نفر و خالی بودن اتاق با خروجی‌های \lr{NotFilled} و \lr{Close} پیاده‌سازی شد.
		\item مسیر \lr{Ent}، \lr{T} و \lr{NotFilled} تنها درخواست‌های مجاز ورود را به فرمان \lr{Open} تبدیل کرد.
		\item سنتز با مصرف بسیار کم منابع و بدون خطا تکمیل شد و فرکانس انتخابی \(100\,\mathrm{MHz}\) پایین‌تر از حد مجاز \(250\,\mathrm{MHz}\) است.
	\end{itemize}

	\section*{نتیجه‌گیری}
	در این آزمایش یک سامانه‌ی ترتیبی کامل با استفاده از طراحی شماتیک ساخته شد. جداسازی شمارنده از منطق کنترل، بررسی مستقل زیرمدار را ممکن کرد و سپس با اتصال آن به آشکارسازهای پر و خالی و حافظه‌ی فرمان در، کنترل ظرفیت و عبور افراد انجام شد. نتایج شبیه‌سازی، تحلیل زمانی و سنتز نرم‌افزاری، عملکرد صحیح طرح را مطابق اهداف آزمایش تأیید می‌کنند.
\end{document}
